\uuid{99Vd}
\titre{Changement de variables}
\auteur{hanani ; mzari}
\organisation{exo7}
\licence{CC BY-NC-SA 4.0}
\datecreate{2020-01-20}
\niveau{L1}
\module{Analyse}
\chapitre{Calcul d'intégrales}
\sousChapitre{Primitives diverses}
\difficulte{2}
\theme{Primitives, changement de variables}
\exo7qcmid{lille-2.primitives.020}

\contenu{
\begin{question}
\qclassification{127.11}
\qtags{motcle=changement de variables}
Parmi les affirmations suivantes, cocher celles qui sont vraies :
\begin{answers}
\bad{Le changement de variable $\displaystyle u=x(1+x^4)$ donne $\displaystyle \int 4x^3(1+x^4)^3\mathrm{d}x=\int 4u^3\mathrm{d}u$.}
\good{Le changement de variable $\displaystyle u=1+x^4$ donne $\displaystyle \int 4x^3(1+x^4)^3\mathrm{d}x=\int u^3\mathrm{d}u$.}
\bad{Le changement de variable $\displaystyle u=2x$ donne $\displaystyle \int \frac{\mathrm{d}x}{1+4x^2}=\int \frac{\mathrm{d}u}{1+u^2}$.}
\good{Le changement de variable $\displaystyle u=2x$ donne $\displaystyle \int \frac{\mathrm{d}x}{1+4x^2}=\frac{1}{2}\int \frac{\mathrm{d}u}{1+u^2}$.}
\end{answers}
\vskip2mm
\begin{explanations}
Avec $u=1+x^4$, on a : $\mathrm{d}u=4x^3\mathrm{d}x$. D'où $\displaystyle \int 4x^3(1+x^4)^3\mathrm{d}x=\int u^3\mathrm{d}u$. De même, avec $u=2x$, on a : $\mathrm{d}u=2\mathrm{d}x$. D'où $\displaystyle \int \frac{\mathrm{d}x}{1+4x^2}=\frac{1}{2}\int \frac{\mathrm{d}u}{1+u^2}$.
\end{explanations}
\end{question}
}
